\section{Framework Architecture}

\subsection{System Overview}

The developed multi-UAV simulation framework is a co-simulation platform for multi-UAV systems that couples flight dynamics, wireless communication and on-board perception withing a single modular architecture. It built on three layers, a flight and coordination layer, a network layer and a perception layer connected through ROS2 which acts as the middleware as shown in Figure. The framework follows four main design principles. First, a layered, modular structure which allows each layer to be developed, replaced or extended. Second, the platform is built entirely using standard open source tools ArduPilot SITL, Gazebo, ROS2 Humble and ns3 which adds the capability of reproducing without proprietary components. YOLOv8n is used to conduct the human detection application. Third, the framework is network aware which it could change the network conditions dynamically rather than applying an idealized channel. Real ROS2/DDS traffic routed through the ns3 simulation, so that modelled propagation and obstacle attenuation act on the actual packets exchanged between nodes. Fourth, the perception stage supports flexible computation, allowing human detection to run either on board each UAV (edge) of at the ground station (ground), which makes the comparison measurable. 

The framework runs on a single Linux workstation (Ubuntu 22.04) on which three UAVs are simulated concurrently. Each UAV combines an ArduPilot SITL instance, a Gazebo iris drone model with a gimble and a camera assembled to the drone model which publishes its telemetry and imagery under an isolated ROS2 namespace (/uav1 - /uav3). A cluster head collects member telemetry providing the needed coordination details for multi-UAV missions. Detection and inter-node traffic traverse through the ns3 channel. To validate the framework against real constrained hardware, the edge configuration is additionally realised in a hardware-in-the-loop (HITL) testbed in which YOLOv8n executes on a physical Raspberry Pi 4B. All other results are obtained on the simulation host.


\subsection{Flight and Coordination Layer}

The simulation framework provides independent flight for each vehicle together with the coordination logic that lets the three UAVs act as a group. Each UAV is driven by a dedicated ArduPilot SITL instance which is coupled to a gazebo model and controlled entirely using ROS2. Since, each UAV topic is published and subscribed under an isolated namespace, the possibility of happening collisions on topics gets avoided. This namespace isolation is the mechanism that makes concurrent multi-UAV simulation possible on one host.

Coordination is built on a clustering architecture. One UAV acts as the cluster head, aggregating the members' telemetry and republishing a consolidated cluster state that reports the head's own flight mode and altitude alongside the real-time positions of the member UAVs. The cluster head holds a primary role with a readily selected backup and a selection mechanism promotes a new head based on selection score based on RSSI and SNR which is combined with a score margin, a consecutive valid condition and a minimum holding time to prevent rapid oscillation. A threshold driven path additionally responds to primary link failure by choosing the best eligible candidate. These controlled handover and failover paths are provided by the architecture but were not tested in the reported experiments.

Multi-UAV behaviour is tested through five phase coordinated missions. All UAVs confirm GPS readiness and wait at synchronization barrier give each UAV its independent time of getting ready for actions such as motor checks and arm checks. Takeoff is done simultaneously to make all three UAVs has the same starting time. The mission executes as simultaneous takeoff, way point navigation and return to launch. Barriers are set in between so the UAVs get synchronized so that no vehicle proceeds before the others are ready which keeps the three vehicles phase aligned throughout the mission.


\subsection{Network Layer}

The network layer is what distinguishes the framework from physics only multi-UAV simulators which assumed ideal link. The network emulates the wireless channel with ns3 and routes the actual inter-node traffic through the channel. \emph{ROS2/DDS} packets are injected into the channel (ns3 simulation) through a TAP bridge, so that the modelled propagation and contention act on the live communication between vehicles rather than on an offline trace. Therefore, a packet carrying a detection experiences the same distance and obstacle dependent conditions the channel model computes and the resulting delay, loss and throughput are properties of the emulated radio which entirely removes the concept of assumptions.

Channel quality is derived per link from two components. A Friis free space propagation loss model provides the baseline distance dependent received power and an obstacle aware loss adds attenuation whenever the straight line path between two endpoints (UAVs) is blocked by the physical obstacles which are placed in the environment determined by ray casting against the obstacles. From these, the model produces per link RSSI and SNR for the full set of links among the three UAVs and the ground station (each UAV-UAV pair and each UAV-GCS pair). The channel state evolves continuously as the vehicles move, since loss depends on separation and line of sight which changes with the vehicle movement. The UAV positions are fed to ns3 using a ros2 topic which consistently publishes each UAVs position through out the entire simulation. The application traffic is mapped onto this channel according to the perception configuration. In the edge processing configuration, inference is performed on board and only the compact detection messages cross the emulated radio path, leaving the imagery on a local, unimpaired path. In the ground processing configuration, compressed imagery is transmitted across the channel for inference at the ground station. This separation of a local data path from the radio path is discussed strongly in the HITL testbed of Section VI.


\subsection{Perception Layer}

This layer is responsible in performing the human detection application using UAV imagery. Detection uses YOLOv8n restricted to the COCO \emph{person} class at a fixed confidence threshold. This model is chosen specifically since it is light weighted and suitable to be used for on-board execution on constrained hardware. Each UAV captures $1280 \times 720$ frames at 5Hz using the downward facing camera assembled to each UAV, simulated by a Gazebo plugin. Same detector and the threshold are used regardless of where inference runs, so that the two configurations differ only in placement, not in the model.

The framework supports two processing configurations that place the detector at opposite ends of the communication link. In the edge processing configuration, inference runs on the UAV's on-board node. Here, only the compact detection messages cross the emulated radio path. In the ground processing configuration, the raw frame is JPEG compressed on the UAV, then transmitted across the ns3 channel and decoded and processed at the ground station. The two configurations therefore trade a small, robust message stream against the transmission of comparatively large imagery over an impaired link, which is precisely the trade off which is quantified in Section V.

To maintain the pipeline real-time when the consumer (processing unit) is slower than the 5Hz source, a best effort QoS profile with a queue depth of one is used by the camera subscriptions. This method ensures that the detector always processes the newest available frame rather than accumulating a backlog of stale ones. This design choice specifically matters because on-board inference on constrained hardware can run slower than the frame rate, discarding stale frames keeps detection latency controlled at the cost of skipping frames than letting delay grow unbounded.
